\documentclass[conference]{IEEEtran}
\IEEEoverridecommandlockouts
\usepackage{cite}
\usepackage{amsmath,amssymb,amsfonts}
\usepackage{algorithmic}
\usepackage{graphicx}
\DeclareGraphicsExtensions{.eps,.pdf,.png}
\usepackage{textcomp}
\usepackage{xcolor}
\usepackage{url}
\usepackage{hyperref}
\usepackage{revhistory}
\newenvironment{raddblock}{\color{red}}{}
\newenvironment{rdelblock}{\color{gray}}{}
\begin{document}

\title{Research Meeting}

\author{\IEEEauthorblockN{Shun Yamachika}
\IEEEauthorblockA{\textit{Department of Informatics, School of Science and Technology} \\
\textit{Kwansei Gakuin University}\\
Sanda, Hyogo, Japan \\
shun@lsnl.jp}}

\maketitle


\section*{現在のストーリー案の骨格}
\label{sec:orga133276}

\begin{itemize}
\item 背景 量子ネットワークにおいて忠実度を高いリンクを効率的に判定する手法 LinkSelFiE が提案されている
\item 動機 LinkSelFiE は通信需要を考慮していないが、現実には通信需要が高くかつ忠実度の高いリンクの判定が望まれる
\item 目的 少ない計測 (バウンス) により利用率 x 忠実度が高いリンクの判定を可能とする
\end{itemize}


\section*{現在の問題設定(先生の対話から抜粋)}
\label{sec:org7674224}

--- 
Q. LinkSelfie の論文の問題を拡張したい。


ノードS〜ノードD 間のリンク集合 Ｌ だけでなく、

ノードS〜ノードD\_n間のリンク集合 L\_n (1 <= n <= N, N はノードSの隣接ノード数)がそれぞれ入力される。

ノードS〜ノードD\_n間の重要度 I\_n (0〜1 の値) が入力として与えられる。



\begin{raddblock}
固定信頼度delta が入力として与えられる。
\end{raddblock}

この時、ノードS〜ノードD\_n のノードペアの中で、

K 個のノードペア (S, D\_s\_1), (S, D\_s\_2), \ldots{} (S, D\_s\_K) における忠実度が最大のリンクをそれぞれ発見する。

\begin{raddblock}
発見するリンクのノードペア数 K は入力として与えられる。
\end{raddblock}

ここで F\_n は (S, D\_n) 間において忠実度の最大のリンクの忠実度である。


\section*{なぜ変更したか}
\label{sec:org2861ca5}

\begin{raddblock}
\begin{itemize}
\item 固定信頼度delta が入力として与えられる。
\end{itemize}
\end{raddblock}

\begin{rdelblock}
\begin{itemize}
\item 総バウンスコスト C が入力として与えられる。
\end{itemize}
\end{rdelblock}


先行研究である、linkselfieの優秀なロジックを活かすことができるため。

\begin{raddblock}
\begin{itemize}
\item 発見するリンクのノードペア数 K は入力として与えられる。
\end{itemize}
\end{raddblock}
\begin{rdelblock}
ただし、発見するリンクのノードペア数 K は、

重要度と忠実度の積の総和、つまり

I = \(\sum\)\_\{k = 1\}\^{}K I\_s\_k * F\_s\_k 

が最大となるように定める。
\end{rdelblock}


重要度と忠実度の積は負にならないので、この変更前の設定だと

発見するリンクのノードペア数 K = 宛先ノードの数

になる。常に全ての宛先を選択することになり、複数の宛先間の重要度を考慮
する意味がなくなるため。







\clearpage
\end{document}